%% kapa: exam=1 round=20 year=2009 period=2 form=A subject=appraisal_law subject_ko=부동산관계법규 seq=6 number=46
%% source: Q-Net 감정평가사 기출문제 (cst00309) artlSeq=5024178 file=extracted/1008/5024178/2009년 제20회 감정평가사 1차시험 기출문제/2교시.hwp
%% conversion: kordoc 4.24.51 -> tools/gen1.py (auto); review=proofread-sonnet; flags=none
\begin{question}{46}
국토의 계획 및 이용에 관한 법령상 토지거래계약의 허가 등에 관하여 판례의 입장이 \uline{아닌} 것은?
\begin{choices}
  \item 토지거래계약에 관한 허가구역으로 지정된 구역 안의 토지에 대하여 거래계약을 체결한 당사자 사이에는 그 계약의 완성을 위하여 서로 협력할 의무가 있다.
  \item 당사자가 허가를 배제하거나 잠탈할 목적으로 허가 없이 계약을 체결한 경우 당해 계약은 확정적 무효이다.
  \item 토지거래계약에 관한 허가구역 내의 토지거래계약의 일방 당사자는 토지거래계약허가 신청절차에 협력하지 않는 상대방에 대하여 그 협력의무의 이행을 청구할 수 있다.
  \item 일방 당사자는 토지거래계약허가 취득절차의 일환으로 당사자 사이의 약정에 기초하여 그 거래의 목적인 토지이용계획 관련 인ㆍ허가절차의 이행을 청구하는 소송을 제기할 수 있다.
  \item 유동적 무효의 상태인 토지거래계약의 일방 당사자는 상대방이 토지거래계약허가 신청절차에 협력하지 아니한다는 사유만으로 그 거래계약 자체를 일방적으로 해제할 수 있다.
\end{choices}
\end{question}
\answer{5}
